\documentclass{article}
\usepackage{graphicx}
\usepackage{amsmath}
\usepackage{amsfonts}
\usepackage{bookmark}
\usepackage{booktabs}
\usepackage{exercises}
\graphicspath{ {./images/} }
\usepackage{indentfirst}
\usepackage{amsthm}
\usepackage{yhmath}
\usepackage{placeins}
\newcommand*{\oldneg}{\mathord{\sim}}
\newcommand\showdiv[1]{\overline{\smash{)}#1}}
\newcommand\Mydiv[1]{%
\smash{\raise.3ex\hbox{$\big)$}}\mkern-8mu
        \overline{\enspace\strut#1}}
\setlength{\tabcolsep}{5pt}
\theoremstyle{definition}
\newtheorem{definition}{Definition}

\setlength{\parindent}{.5in}
\setlength{\parskip}{12pt}
\renewcommand{\arraystretch}{2}
\newcommand{\?}{\stackrel{?}{=}}

\title{Turing Paper Notes}
\begin{document}
\maketitle

\setion*{Section 2}
What is meant when by a number differing by an integer from the number computed by a circle-free machine (\texit{Computable sequences and numbers} section)?

Is it assumed that, after going to a new square on the tape, that square contains no symbol unless otherwise specified?

In the example of the machine that computes the sequence $010101...$, why, after, the configuration determined by the $m$-config $b$ and no symbol, is 
$P0,R$ done, followed by going to the next configuration, which also does $R$. This seems like it is shifting over to places to the right, 
and thus missing one place on the tape (this same question comes up with the third and fourth configurations).

For the second example machine, what is the first configuration, is it the $m$-config $b$ with any symbol or no symbol?

We should go over the second example machine, as I was having trouble understanding how it really computed the desired sequence.

What exactly did he mean when he was talking about writing figures on "alternate squares"? Does this have something to do with what he said before: "We also arrange that in the 
sequence of figures on alternate squares there shall be no blanks."
